\subsection{数列と漸化式}\label{節：数列と漸化式}
数を順に並べたものを\bookterm{sequence}{数列}という.
数列のうち,要素が有限個の数列を有限数列,無限個のものを無限数列と呼ぶ.
また,数列の各要素は法則に従っていても,ランダムに並べられていても構わない.
本書では一定の規則(\bookterm{recurrence}{漸化式})によって数が変化する無限数列について扱う.
今後数列と言われた場合には,無限数列を指す.

並べられた各要素を項といい,各項を$a_1,a_2,a_3,\ldots$と書き,添字でそれぞれの順番を表す.

例えば,$a_1=1$から始めて,直前の項に$2$足し続けていくような数列を考えると,
\[
  1,\ 3,\ 5,\ 7,\ \ldots
\]
という\bookterm{sequence}{数列}ができる.この規則を式で書けば,
\[
  a_{n+1}=a_n+2\qquad(n\ge1)
\]
となる.
また,同時に,
\[
  a_n=2n-1 \qquad (n\ge1)
\]
とも表すことができる.

前者のように,項どうしの関係によって\bookterm{sequence}{数列}を定める式を,\bookterm{recurrence}{漸化式}という.
最初に与える値は\bookterm{initial}{初期条件}と呼ばれる.
後者は数列の一般項を表す.

漸化式に関して,先の例では2項間の関係を表していたので,先程の例は
\[
  a_{n}=a_{n-1}+2\qquad(n\ge1)
\]
のようにも,
\[
  a_{n+100}=a_{n+99}+2\qquad(n\ge1)
\]
のようにも書くことができる.
ただし,引き算をする場合は添字が$n$が定義域のうちに存在するかを確認しなればならない.

次の項が一意に計算できる限り,同じ\bookterm{initial}{初期条件}と規則から同じ\bookterm{sequence}{数列}ができる.
また,フィボナッチ数列
\begin{align*}
  &a_1=1,\quad a_2=1\\
  &a_{n+2}=a_{n+1}+a_n
\end{align*}
のように,直前の二項を使う規則では,通常,初めの二項を与える.

本書では漸化式を一般項に変形するためのテクニックや,その裏に隠れる数学的事実を扱う.